\documentclass[11pt,twoside,openright]{book}
\usepackage{springer_monograph}
\usepackage{placeins}
\usepackage{float}

\hypersetup{
    pdftitle={Chapter 02: Mass Spectrometry Principles and Instrument Physics},
    pdfauthor={Lucas VHH TRAN}
}

\begin{document}

% --- Chapter Title Block ---
\setcounter{chapter}{1}
\chapter{Mass Spectrometry Principles and Instrument Physics}
\label{chap:what_is_ms}

\noindent
\textbf{Lucas VHH TRAN}\footnote{CK-CARE Davos, Christine K\"uhne--Center for Allergy Research and Education, Davos, Switzerland. Email: \texttt{lucas.tran@ck-care.ch}}\\
\rule{\textwidth}{0.8pt}

\begin{abstract}
\noindent
Mass spectrometers operate as physical sensors that convert gas-phase ions into digitized mass-to-charge ratios ($m/z$) and ion-current intensities. Neutral molecules cannot be detected directly; they must acquire electric charge through electron ionization (EI), electrospray ionization (ESI) \cite{fenn1989electrospray}, or matrix-assisted laser desorption/ionization (MALDI) \cite{karas1988laser}. Because the instrument measures $m/z$ rather than uncharged mass, molecular identity must be derived by deconvoluting integer charge states ($z$) from isotopic peak spacing. Instrument mass analyzers---quadrupoles \cite{paul1990electromagnetic}, time-of-flight (TOF) \cite{guilhaus1995principles}, and Orbitraps \cite{makarov2000orbitrap, marshall1998fourier}---differ by orders of magnitude in resolving power and mass accuracy \cite{zubarev2007persistent, gross2017mass}, setting physical boundaries on isobaric separation. Furthermore, the choice of acquisition strategy governs data completeness: data-dependent acquisition (DDA) introduces stochastic missingness through top-$N$ precursor selection \cite{domon2010options, michalski2011more}, whereas data-independent acquisition (DIA) captures multiplexed fragmentation arrays at the expense of computational deconvolution \cite{gillet2012targeted, venable2004automated, demichev2020dia}. This chapter presents the physical governing equations, charge deconvolution algorithms, resolving power thresholds, and matrix suppression kinetics that define mass spectrometry data structures in R 4.6.1 \cite{domon2006ms, gross2017mass}.
\end{abstract}

\vspace{1em}
\noindent\textbf{Keywords:}\enspace Mass-to-charge ratio, Charge state deconvolution, Resolving power, Isobaric separation, DDA vs DIA, Ion suppression, Stable isotope standards.

\vspace{1em}
\section*{Chapter Overview and Architectural Roadmap}

Table~\ref{tab:ms_physics_roadmap} outlines the physical principles and analytical consequences of mass spectrometry instrumentation.

\begin{table}[htbp]
\centering
\caption{Physical principles and data characteristics of mass spectrometry systems.}
\label{tab:ms_physics_roadmap}
\small
\begin{tabularx}{\textwidth}{>{\raggedright\arraybackslash}l >{\raggedright\arraybackslash}p{4.2cm} >{\raggedright\arraybackslash}X >{\raggedright\arraybackslash}p{2.8cm}}
\toprule
\textbf{Physical Domain} & \textbf{Instrument Mechanism} & \textbf{Mathematical Formulation} & \textbf{Data Structure Impact} \\
\midrule
\textbf{1. Ionization} & Gas-phase charging (ESI/MALDI) \cite{fenn1989electrospray, karas1988laser} & $m/z = (M + z \cdot m_{\text{proton}})/z$ & Multi-charge envelopes \\
\addlinespace
\textbf{2. Analyzer} & Trajectory dispersion in fields \cite{makarov2000orbitrap, guilhaus1995principles} & $R = m / \Delta m_{\text{FWHM}}$ \cite{zubarev2007persistent} & Peak width and ppm error \\
\addlinespace
\textbf{3. Acquisition} & DDA stochastic precursor selection \cite{michalski2011more} & Top-$N$ instantaneous intensity rank \cite{domon2010options} & Stochastic missingness \\
\addlinespace
\textbf{4. Multiplexing} & DIA windowed fragmentation \cite{gillet2012targeted} & Swath stepping $\Delta m = 20$--$25$ Da & Chimeric fragment arrays \cite{demichev2020dia} \\
\addlinespace
\textbf{5. Matrix Effect} & Droplet evaporation competition \cite{king2000mechanisms} & $I_{\text{obs}} = S(t) \cdot k \cdot C_{\text{analyte}}$ & Ion suppression loss \cite{annesley2003ion} \\
\bottomrule
\end{tabularx}
\end{table}

\begin{notebox}
\textbf{Tutor Warning: Mass-to-Charge is Not Molecular Mass}\\[0.4em]
A mass spectrometer measures the dimensionless ratio $m/z$, not molecular mass. A single peak observed at $m/z = 450.25$ represents a 449.24 Da singly-charged ion ($z=1$), an 898.49 Da doubly-charged peptide ($z=2$), or a 1,347.73 Da triply-charged species ($z=3$). Attempting database matching without resolving the charge state from carbon-13 isotopic spacing ($\Delta(m/z) \approx 1.00335/z$) yields false identifications.
\end{notebox}

% =============================================================================
\section{The Physical Measurement: Mass-to-Charge Ratio}
\label{sec:mz_measurement}
% =============================================================================

Electric and magnetic fields manipulate charged particles according to the Lorentz force \cite{gross2017mass, paul1990electromagnetic}:
\begin{equation}
\mathbf{F} = q(\mathbf{E} + \mathbf{v} \times \mathbf{B})
\end{equation}
Because acceleration $\mathbf{a} = \mathbf{F}/m$ depends directly on the ratio of charge to inertial mass ($q/m$), mass spectrometers separate ions by their mass-to-charge ratio ($m/z$) \cite{domon2006ms}.

\subsection{Charge State Determination from Carbon-13 Isotope Spacing}

In organic molecules, naturally occurring carbon contains approximately $98.93\%$ $^{12}\text{C}$ and $1.07\%$ $^{13}\text{C}$ \cite{gross2017mass}. The mass difference between a $^{13}\text{C}$ atom and a $^{12}\text{C}$ atom is:
\begin{equation}
\Delta m_{^{13}\text{C} - ^{12}\text{C}} = 1.003355\;\text{Da}
\end{equation}
When an ion acquires integer charge $z$, the observed difference in $m/z$ between the monoisotopic peak ($M_0$) and the first isotopic peak ($M_1$) shrinks in inverse proportion to $z$:
\begin{equation}
\Delta(m/z) = \frac{1.003355}{z} \implies z = \text{round}\left( \frac{1.003355}{\Delta(m/z)} \right)
\end{equation}
Once the charge state $z$ is inferred, the neutral monoisotopic mass $M_{\text{neutral}}$ is reconstructed by subtracting the mass of the charge-carrying protons ($m_{\text{proton}} = 1.007276$ Da):
\begin{equation}
M_{\text{neutral}} = z \cdot (m/z)_{\text{obs}} - z \cdot m_{\text{proton}}
\end{equation}

% =============================================================================
\section{Mass Resolving Power and Mass Accuracy}
\label{sec:resolving_power}
% =============================================================================

Two orthogonal metrics define mass analyzer performance: resolving power ($R$) and mass measurement accuracy ($\delta_{\text{ppm}}$) \cite{zubarev2007persistent, makarov2000orbitrap, marshall1998fourier}.

\subsection{Resolving Power Definition}

The mass resolving power $R$ describes the ability of an analyzer to separate adjacent spectral peaks of equal intensity. Under the Full Width at Half Maximum (FWHM) convention \cite{zubarev2007persistent}:
\begin{equation}
R = \frac{m}{\Delta m_{\text{FWHM}}} \implies \Delta m_{\text{FWHM}} = \frac{m}{R}
\end{equation}
To separate two isobaric compounds with mass difference $\Delta m = |m_2 - m_1|$, the instrument resolving power must satisfy:
\begin{equation}
R_{\text{required}} \ge \frac{m_1}{\Delta m}
\end{equation}

\subsection{Mass Accuracy in Parts-Per-Million}

Mass measurement error represents the deviation between the observed experimental peak apex and the theoretical elemental mass:
\begin{equation}
\delta_{\text{ppm}} = \frac{m_{\text{obs}} - m_{\text{theor}}}{m_{\text{theor}}} \times 10^6
\end{equation}
High-resolution analyzers operate at tolerances $|\delta_{\text{ppm}}| \le 5.0$ ppm, enabling narrow extracted ion chromatogram (EIC) windows that exclude background chemical noise.

% =============================================================================
\section{Acquisition Architectures: DDA versus DIA}
\label{sec:dda_vs_dia}
% =============================================================================

Untargeted mass spectrometry operates under two distinct tandem acquisition regimes \cite{domon2010options}:

\subsection{Data-Dependent Acquisition (DDA)}

In DDA, each duty cycle begins with an MS1 survey scan \cite{michalski2011more}. The instrument firmware identifies the $N$ most abundant precursor peaks (Top-$N$ protocol), isolates each ion sequentially using a narrow quadrupole transmission window ($\approx 1.0$--$1.5$ Da), fragments it via collision-induced dissociation (CID) or higher-energy collisional dissociation (HCD), and records clean, single-precursor fragment spectra.

\noindent\textbf{Analytical Implication:}\enspace Because complex proteomes contain hundreds of co-eluting peptides per second, Top-$N$ selection is stochastic. Lower-abundance precursors hover near the detection threshold; Poisson ion-arrival statistics determine whether a peptide is fragmented in a given run, creating stochastic missingness across cohort runs.

\subsection{Data-Independent Acquisition (DIA)}

In DIA, the instrument disregards instantaneous precursor intensities \cite{gillet2012targeted, venable2004automated}. Instead, it systematically steps through a series of wide precursor isolation windows ($\Delta m = 20$--$25$ Da) spanning the entire mass range. All ions within each window are fragmented simultaneously.

\noindent\textbf{Analytical Implication:}\enspace DIA provides $100\%$ acquisition completeness across all runs. However, the resulting MS2 spectra are multiplexed mixtures of fragments from dozens of co-isolated precursors. Identification requires matched spectral libraries and algorithmic deconvolution (e.g., DIA-NN \cite{demichev2020dia}, Spectronaut).

% =============================================================================
\section{Matrix Effects and Ion Suppression}
\label{sec:ion_suppression}
% =============================================================================

In electrospray ionization (ESI), charged droplets formed at atmospheric pressure undergo solvent evaporation and Coulomb fission \cite{fenn1989electrospray, king2000mechanisms}. Co-eluting non-volatile salts, phospholipids, or polymer contaminants compete with target analytes for droplet surface access and available charge.

This competition attenuates analyte ionization by an empirical suppression factor $S(t) \in [0, 1]$ \cite{matuszewski2003strategies, annesley2003ion}:
\begin{equation}
I_{\text{obs}}(t) = S(t) \cdot k \cdot C_{\text{analyte}}
\end{equation}
Severe suppression causes signals to drop below the limit of detection, generating Missing Not At Random (MNAR) patterns.

\noindent\textbf{Correction via Stable Isotope Standards:}\enspace
Spiking a stable isotope-labeled (SIL) internal standard corrects this distortion \cite{ong2002stable}. Because the SIL analog co-elutes with the native analyte, both experience the identical suppression factor $S(t)$:
\begin{equation}
\frac{I_{\text{analyte}}}{I_{\text{SIL}}} = \frac{S(t) \cdot k \cdot C_{\text{analyte}}}{S(t) \cdot k \cdot C_{\text{SIL}}} = \frac{C_{\text{analyte}}}{C_{\text{SIL}}}
\end{equation}
The matrix suppression factor cancels algebraically, preserving accurate linear quantification.

% =============================================================================
\section{Laboratory Workflow and Executable Implementations}
\label{sec:lab_implementation}
% =============================================================================

The laboratory exercises demonstrate charge deconvolution from isotope clusters, resolving power modeling across mass analyzers, DDA stochastic missingness simulation, and ESI ion suppression kinetics in R 4.6.1.

\subsection{Listing 2.1: Resolving Charge States from Isotope Spacing}

\begin{rchunk}{Listing 2.1: Deconvoluting Charge States from Isotopic Deltas}
library(tibble)
library(dplyr)
library(tidyr)

m_proton <- 1.007276466879
c13_diff <- 1.003354835

# Simulated isotopic clusters: z = 1, 2, 3, 20
test_charges <- c(1L, 2L, 3L, 20L)
mz_monos     <- c(450.2500, 450.2500, 450.2500, 848.5824)

deconv_list <- list()
for (i in seq_along(test_charges)) {
  z_true   <- test_charges[i]
  mz_0     <- mz_monos[i]
  delta_mz <- c13_diff / z_true
  mz_1     <- mz_0 + delta_mz
  
  meas_delta    <- round(mz_1 - mz_0, 4)
  inferred_z    <- round(c13_diff / meas_delta)
  neutral_mass  <- inferred_z * mz_0 - inferred_z * m_proton
  
  deconv_list[[i]] <- tibble(
    True_Z       = z_true,
    Observed_MZ0 = mz_0,
    Delta_MZ     = meas_delta,
    Inferred_Z   = as.integer(inferred_z),
    Neutral_Mass = round(neutral_mass, 2)
  )
}

print(as.data.frame(bind_rows(deconv_list)))
\end{rchunk}

\begin{routput}
  True\_Z Observed\_MZ0 Delta\_MZ Inferred\_Z Neutral\_Mass
1      1      450.2500   1.0034          1       449.24
2      2      450.2500   0.5017          2       898.49
3      3      450.2500   0.3345          3      1347.73
4     20      848.5824   0.0502         20     16951.50
\end{routput}

\subsection{Listing 2.2: Resolving Power and Isobaric Separation}

\begin{rchunk}{Listing 2.2: Calculating Peak FWHM and Isobaric Thresholds}
analyzers <- tibble(
  Analyzer  = c("Unit Quadrupole", "Time-of-Flight", "Orbitrap Standard", "Orbitrap Ultra"),
  Nominal_R = c(1000, 30000, 60000, 240000)
) %>%
  mutate(FWHM_at_500 = round(500.0 / Nominal_R, 5))

print(as.data.frame(analyzers))

# Isobaric doublet: Glutamine (146.0691) vs Lysine (146.1055)
delta_iso <- abs(146.1055 - 146.0691)
req_R     <- round(146.0691 / delta_iso)
cat(sprintf("\nGlutamine vs Lysine Delta: %.4f Da | Required R >= %d\n", delta_iso, req_R))
\end{rchunk}

\begin{routput}
           Analyzer Nominal\_R FWHM\_at\_500
1   Unit Quadrupole      1000     0.50000
2    Time-of-Flight     30000     0.01667
3 Orbitrap Standard     60000     0.00833
4    Orbitrap Ultra    240000     0.00208

Glutamine vs Lysine Delta: 0.0364 Da | Required R >= 4013
\end{routput}

\subsection{Listing 2.3: Simulating Stochastic Sampling in DDA}

\begin{rchunk}{Listing 2.3: Quantifying Missingness Driven by TopN Selection}
set.seed(42)
n_prec <- 300; n_reps <- 5
true_abund <- rlnorm(n_prec, meanlog = 12, sdlog = 2.0)

dda_mat <- matrix(FALSE, nrow = n_prec, ncol = n_reps)
for (rep_i in 1:n_reps) {
  obs_int <- true_abund * rlnorm(n_prec, 0, 0.15)
  dda_mat[order(obs_int, decreasing = TRUE)[1:25], rep_i] <- TRUE
}

cat(sprintf("Precursors Detected in All 5 Replicates: %.1f%%\n",
            mean(rowSums(dda_mat) == n_reps) * 100))
cat(sprintf("Precursors Completely Missing (0 of 5):   %.1f%%\n",
            mean(rowSums(dda_mat) == 0) * 100))
\end{rchunk}

\begin{routput}
Precursors Detected in All 5 Replicates: 5.3%
Precursors Completely Missing (0 of 5):   89.7%
\end{routput}

\subsection{Listing 2.4: ESI Ion Suppression Kinetics}

\begin{rchunk}{Listing 2.4: Calibration Curve Slope Attenuation Under Matrix}
conc <- c(1, 5, 10, 25, 50, 100)
resp_clean  <- 10000 * conc + rnorm(length(conc), 0, 5000)
resp_matrix <- 10000 * 0.35 * conc + rnorm(length(conc), 0, 3000)

fit_c <- lm(resp_clean ~ conc)
fit_m <- lm(resp_matrix ~ conc)

cat(sprintf("Clean Slope:  %.0f counts/nM (R2 = %.4f)\n", coef(fit_c)[2], summary(fit_c)$r.squared))
cat(sprintf("Matrix Slope: %.0f counts/nM (R2 = %.4f)\n", coef(fit_m)[2], summary(fit_m)$r.squared))
cat(sprintf("Signal Loss:  %.1f%%\n", (1 - coef(fit_m)[2] / coef(fit_c)[2]) * 100))
\end{rchunk}

\begin{routput}
Clean Slope:  10012 counts/nM (R2 = 0.9999)
Matrix Slope: 3461 counts/nM (R2 = 0.9996)
Signal Loss:  65.4%
\end{routput}

% =============================================================================
\section{Five Deadly Sins of Mass Spectrometry Data Interpretation}
\label{sec:deadly_sins}
% =============================================================================

\begin{enumerate}
    \item \textbf{Confusing $m/z$ with Molecular Mass:}\enspace Assuming observed peaks represent singly-charged ions misidentifies multiply-charged peptides and proteins.
    \item \textbf{Assuming DDA Completeness:}\enspace Treating missing values in DDA as technical errors rather than stochastic precursor sampling distorts differential abundance statistics.
    \item \textbf{Overlooking Matrix Suppression:}\enspace Failing to monitor internal standards in complex matrices attributes technical ionization losses to biological downregulation.
    \item \textbf{Neglecting Isobaric Interferences:}\enspace Using broad extraction windows on low-resolution analyzers merges distinct chemical species into single distorted peaks.
    \item \textbf{Ignoring In-Source Fragmentation:}\enspace Treating source-generated fragments as biological metabolites inflates feature tables with artificial molecular entities.
\end{enumerate}

% =============================================================================
\section{Pre-Flight Diagnostic Checklist}
\label{sec:diagnostic_checklist}
% =============================================================================

Table~\ref{tab:ms_physics_checklist} provides an audit checklist to verify instrument parameters prior to data interpretation.

\begin{table}[htbp]
\centering
\caption{Diagnostic checklist for mass spectrometry instrument configurations.}
\label{tab:ms_physics_checklist}
\small
\begin{tabularx}{\textwidth}{>{\raggedright\arraybackslash}l >{\raggedright\arraybackslash}p{5.2cm} >{\raggedright\arraybackslash}X c}
\toprule
\textbf{Category} & \textbf{Diagnostic Check} & \textbf{Acceptance Criteria} & \textbf{Status} \\
\midrule
\textbf{Charge} & Isotope cluster spacing verified. & Peak spacing $\Delta(m/z) \approx 1.00335/z$. & [ ] \\
\addlinespace
\textbf{Resolution} & Resolving power adequate for target isobars. & $R \ge m / \Delta m$ for critical metabolite pairs. & [ ] \\
\addlinespace
\textbf{Accuracy} & Mass calibration verified against lock mass. & Absolute error $|\delta_{\text{ppm}}| \le 5.0$ ppm. & [ ] \\
\addlinespace
\textbf{Strategy} & Acquisition mode matches statistical model. & DDA modeled with MNAR/MAR; DIA deconvoluted. & [ ] \\
\addlinespace
\textbf{Suppression} & Stable isotope labeled standards tracked. & Internal standard CV $\le 15\%$ across all runs. & [ ] \\
\bottomrule
\end{tabularx}
\end{table}

\FloatBarrier

% =============================================================================
\section{Chapter Exercises}
\label{sec:chapter_exercises}
% =============================================================================

\noindent\textbf{Exercise 2.1 (Charge State and Neutral Mass Reconstruction).}\enspace
\label{ex:charge_deconvolution}
An unknown tryptic peptide is detected on an LC-MS system. The monoisotopic peak appears at $m/z = 624.8348$, and the first $^{13}\text{C}$ isotopic peak appears at $m/z = 625.1693$.
\begin{enumerate}
    \item Determine the peak spacing $\Delta(m/z)$ and calculate the integer charge state $z$.
    \item Reconstruct the monoisotopic neutral mass ($M_{\text{neutral}}$) in Daltons.
\end{enumerate}

\vspace{0.8em}
\noindent\textbf{Exercise 2.2 (Resolving Power for Isobaric Metabolites).}\enspace
\label{ex:resolving_power_isobars}
Two small molecules have identical nominal mass ($132$ Da) but distinct elemental compositions: Metabolite 1 ($132.0773$ Da) and Metabolite 2 ($132.1025$ Da).
\begin{enumerate}
    \item Calculate the exact mass difference $\Delta m$.
    \item Compute the minimum mass resolving power $R$ required for baseline separation and state which mass analyzers satisfy this requirement.
\end{enumerate}

\vspace{0.8em}
\noindent\textbf{Exercise 2.3 (Stochastic Missingness Mechanisms in DDA).}\enspace
\label{ex:dda_stochastic_missingness}
Explain why Data-Dependent Acquisition (DDA) produces stochastic missing values across replicate runs of the same biological sample. Categorize this missingness under Rubin's missing data taxonomy (MCAR, MAR, MNAR).

\vspace{0.8em}
\noindent\textbf{Exercise 2.4 (Matrix Suppression Correction via Isotope Dilution).}\enspace
\label{ex:sil_matrix_correction}
A targeted LC-MS assay suffers $60\%$ electrospray ion suppression ($S(t) = 0.40$) during the elution window of an endogenous analyte.
\begin{enumerate}
    \item Write the mathematical response equations for the endogenous analyte and a co-eluting stable isotope-labeled (SIL) internal standard.
    \item Prove that the area ratio $I_{\text{analyte}} / I_{\text{SIL}}$ is invariant to the suppression factor.
\end{enumerate}

\FloatBarrier

% =============================================================================
\section{Worked Solutions to Chapter Exercises}
\label{sec:solutions}
% =============================================================================

\subsection{Solution to Exercise 2.1: Charge State and Neutral Mass Reconstruction}
\noindent\textbf{Part 1: Charge State Determination}\enspace
The observed isotopic peak spacing is:
\begin{equation}
\Delta(m/z) = 625.1693 - 624.8348 = 0.3345\;\text{Da}
\end{equation}
The integer charge state is:
\begin{equation}
z = \text{round}\left( \frac{1.003355}{0.3345} \right) = \text{round}(2.9996) = 3
\end{equation}

\noindent\textbf{Part 2: Neutral Mass Calculation}\enspace
Subtracting 3 protons ($m_{\text{proton}} = 1.007276$ Da):
\begin{equation}
M_{\text{neutral}} = 3 \times 624.8348 - 3 \times 1.007276 = 1874.5044 - 3.0218 = 1871.4826\;\text{Da}
\end{equation}

\subsection{Solution to Exercise 2.2: Resolving Power for Isobaric Metabolites}
\noindent\textbf{Part 1: Mass Difference}\enspace
\begin{equation}
\Delta m = 132.1025 - 132.0773 = 0.0252\;\text{Da} = 25.2\;\text{mDa}
\end{equation}

\noindent\textbf{Part 2: Resolving Power Requirement}\enspace
\begin{equation}
R \ge \frac{132.0773}{0.0252} \approx 5,241
\end{equation}
A unit-mass quadrupole ($R \approx 1,000$) fails to resolve the doublet, merging them into a single envelope. A Time-of-Flight ($R \ge 30,000$) or Orbitrap ($R \ge 60,000$) resolves the doublet completely.

\subsection{Solution to Exercise 2.3: Stochastic Missingness Mechanisms in DDA}
In DDA, the mass spectrometer dynamically selects only the top $N$ most abundant precursor ions detected in each survey scan. Because low-abundance peptides fluctuate near the selection threshold, small run-to-run variations in ionization or chromatographic apex timing determine whether a peptide is selected or excluded. Under Rubin's taxonomy:
\begin{itemize}
    \item Peptides below the instrument detection limit exhibit Missing Not At Random (MNAR) censoring.
    \item Peptides with sufficient abundance that fluctuate across the Top-$N$ boundary exhibit Missing at Random (MAR) stochasticity conditioned on instantaneous MS1 intensity.
\end{itemize}

\subsection{Solution to Exercise 2.4: Matrix Suppression Correction via Isotope Dilution}
\noindent\textbf{Part 1: Response Equations}\enspace
Let $k$ denote the instrument response factor. Under $60\%$ suppression ($S = 0.40$):
\begin{align}
I_{\text{analyte}} &= 0.40 \cdot k \cdot C_{\text{analyte}} \\
I_{\text{SIL}}     &= 0.40 \cdot k \cdot C_{\text{SIL}}
\end{align}

\noindent\textbf{Part 2: Invariance Proof}\enspace
Computing the signal ratio:
\begin{equation}
\frac{I_{\text{analyte}}}{I_{\text{SIL}}} = \frac{0.40 \cdot k \cdot C_{\text{analyte}}}{0.40 \cdot k \cdot C_{\text{SIL}}} = \frac{C_{\text{analyte}}}{C_{\text{SIL}}}
\end{equation}
The suppression coefficient $0.40$ and instrument factor $k$ cancel algebraically, preserving linear calibration accuracy.

% =============================================================================
\begin{thebibliography}{99}

\bibitem{fenn1989electrospray}
Fenn, J.B., Mann, M., Meng, C.K., Wong, S.F., Whitehouse, C.M.:
\newblock Electrospray ionization for mass spectrometry of large biomolecules.
\newblock \emph{Science} \textbf{246}(4926), 64--71 (1989).
\newblock \doi{10.1126/science.246.4926.64}

\bibitem{karas1988laser}
Karas, M., Hillenkamp, F.:
\newblock Laser desorption ionization of proteins with molecular masses exceeding 10,000 daltons.
\newblock \emph{Analytical Chemistry} \textbf{60}(20), 2299--2301 (1988).
\newblock \doi{10.1021/ac00171a028}

\bibitem{gross2017mass}
Gross, J.H.:
\newblock Mass Spectrometry: A Textbook, 3rd edn.
\newblock Springer, Cham (2017).
\newblock \doi{10.1007/978-3-319-54398-7}

\bibitem{domon2006ms}
Domon, B., Aebersold, R.:
\newblock Mass spectrometry and protein analysis.
\newblock \emph{Science} \textbf{312}(5771), 212--217 (2006).
\newblock \doi{10.1126/science.1124619}

\bibitem{makarov2000orbitrap}
Makarov, A.:
\newblock Electrostatic axially harmonic orbital trapping: a high-performance technique of mass analysis.
\newblock \emph{Analytical Chemistry} \textbf{72}(6), 1156--1162 (2000).
\newblock \doi{10.1021/ac991131p}

\bibitem{marshall1998fourier}
Marshall, A.G., Hendrickson, C.L., Jackson, G.S.:
\newblock Fourier transform ion cyclotron resonance mass spectrometry: a primer.
\newblock \emph{Mass Spectrometry Reviews} \textbf{17}(1), 1--35 (1998).
\newblock \doi{10.1002/(SICI)1098-2787(1998)17:1<1::AID-MAS1>3.0.CO;2-K}

\bibitem{paul1990electromagnetic}
Paul, W.:
\newblock Electromagnetic traps for charged particles.
\newblock \emph{Reviews of Modern Physics} \textbf{62}(3), 531--540 (1990).
\newblock \doi{10.1103/RevModPhys.62.531}

\bibitem{guilhaus1995principles}
Guilhaus, M.:
\newblock Principles and instrumentation in time-of-flight mass spectrometry. Special feature: tutorial.
\newblock \emph{Journal of Mass Spectrometry} \textbf{30}(11), 1519--1532 (1995).
\newblock \doi{10.1002/jms.1190301102}

\bibitem{zubarev2007persistent}
Zubarev, R., Mann, M.:
\newblock On the proper use of the terms mass resolution and mass resolving power.
\newblock \emph{Analytical Chemistry} \textbf{79}(20), 7937--7937 (2007).
\newblock \doi{10.1021/ac072023d}

\bibitem{domon2010options}
Domon, B., Aebersold, R.:
\newblock Options and considerations when selecting a quantitative proteomics strategy.
\newblock \emph{Nature Biotechnology} \textbf{28}(7), 710--721 (2010).
\newblock \doi{10.1038/nbt.1661}

\bibitem{michalski2011more}
Michalski, A., Cox, J., Mann, M.:
\newblock More than 100,000 detectable peptide features and more than 10,000 identified proteins in a single analysis of the human proteome.
\newblock \emph{Journal of Proteome Research} \textbf{10}(4), 1785--1793 (2011).
\newblock \doi{10.1021/pr101065u}

\bibitem{gillet2012targeted}
Gillet, L.C., Navarro, P., Tate, S., et~al.:
\newblock Targeted data extraction of the {MS/MS} spectra generated by data-independent acquisition: a new concept for consistent and accurate proteome analysis.
\newblock \emph{Molecular \& Cellular Proteomics} \textbf{11}(6), O111.016717 (2012).
\newblock \doi{10.1074/mcp.O111.016717}

\bibitem{venable2004automated}
Venable, J.D., Dong, M.Q., Wohlschlegel, J., Dillin, A., Yates, J.R.:
\newblock Automated approach for quantitative analysis of complex peptide mixtures from tandem mass spectra.
\newblock \emph{Nature Methods} \textbf{1}(1), 39--45 (2004).
\newblock \doi{10.1038/nmeth705}

\bibitem{demichev2020dia}
Demichev, V., Messner, C.B., Vernardis, S.I., Lilley, K.S., Ralser, M.:
\newblock {DIA-NN}: neural networks and interference correction enable deep proteome coverage in high throughput.
\newblock \emph{Nature Methods} \textbf{17}(1), 41--44 (2020).
\newblock \doi{10.1038/s41592-019-0638-x}

\bibitem{king2000mechanisms}
King, R., Bonfiglio, R., Fernandez-Metzler, C., Miller-Stein, C., Olah, T.:
\newblock Mechanisms of matrix effects in electrospray ionization {LC-MS}.
\newblock \emph{Journal of the American Society for Mass Spectrometry} \textbf{11}(11), 942--950 (2000).
\newblock \doi{10.1016/S1044-0305(00)00163-X}

\bibitem{matuszewski2003strategies}
Matuszewski, B.K., Constanzer, M.L., Chavez-Eng, C.M.:
\newblock Strategies for the assessment of matrix effect in quantitative bioanalytical methods based on {HPLC-MS/MS}.
\newblock \emph{Analytical Chemistry} \textbf{75}(13), 3019--3030 (2003).
\newblock \doi{10.1021/ac020361s}

\bibitem{annesley2003ion}
Annesley, T.M.:
\newblock Ion suppression in mass spectrometry.
\newblock \emph{Clinical Chemistry} \textbf{49}(7), 1041--1044 (2003).
\newblock \doi{10.1373/49.7.1041}

\bibitem{ong2002stable}
Ong, S.E., Blagoev, B., Kratchmarova, I., et~al.:
\newblock Stable isotope labeling by amino acids in cell culture, {SILAC}, as a simple and accurate approach to expression proteomics.
\newblock \emph{Molecular \& Cellular Proteomics} \textbf{1}(5), 376--386 (2002).
\newblock \doi{10.1074/mcp.M200025-MCP200}

\end{thebibliography}

\end{document}
